\documentclass[12pt,landscape]{article}
\usepackage{amsmath}
\usepackage{pdflscape}
\usepackage{color}
\usepackage[margin=1in]{geometry}

\usepackage{fancyhdr}
\usepackage{totpages}

\fancyhead[L]{Purdue University, ME 50900}
\fancyhead[R]{Prof.\ Ivan C.\ Christov}
\fancyfoot[C]{Page \thepage\ of \ref*{TotPages}}

\usepackage{accents}

\DeclareMathOperator{\tr}{tr}

\usepackage[most]{tcolorbox}

\newtcolorbox{empheqboxed}{colback=gray!5,
 colframe=black,
 width=\textwidth,
 sharpish corners,
 top=-2mm, % default value 2mm
}

\begin{document}

\thispagestyle{fancy}

\Large

\begin{minipage}[t]{0.45\textwidth}
\centerline{\fbox{Compressible flow}}
\begin{align}
\frac{\partial \rho}{\partial t} + \nabla\cdot(\rho \underline{v}) &= 0 \tag*{\text{\small[COM v1]}}\\
\partial_t\rho + \partial_j (\rho v_j) &= 0 \nonumber
\end{align}
\textbf{\emph{equivalently}}
\[ \frac{D \rho}{D t} + \rho\nabla\cdot\underline{v} = 0 \eqno\text{\small[COM v2]}\]
\[ \partial_t\rho + v_j\partial_j \rho + \rho \partial_j v_j = 0 \]

\noindent\makebox[\linewidth]{\rule{\textwidth}{1pt}}

\vspace{-3mm}

\[ \frac{\partial (\rho \underline{v})}{\partial t} + \nabla\cdot(\rho \underline{v} \, \underline{v}) = \rho\underline{F} + \nabla\cdot\underline{\underline{T}}
\eqno\text{\small[COLM v1]}\]
\[ \partial_t (\rho v_i) + \partial_j(\rho v_j v_i) = \rho F_i + \partial_j T_{ji} \]

\textbf{\emph{equivalently}} (using COM to simplify)
\[ \rho\frac{D \underline{v}}{D t} = \rho\underline{F} + \nabla\cdot\underline{\underline{T}} \eqno\text{\small[COLM v2]}\]
\[\rho \left(\partial_t v_i + v_j\partial_j v_i\right) = \rho F_i + \partial_j T_{ji}\]

\noindent\makebox[\linewidth]{\rule{\textwidth}{1pt}}

\end{minipage}
\hfill
\begin{minipage}[t]{0.45\textwidth}
\centerline{\fbox{Incompressible flow}}

\begin{align}
\nabla\cdot \underline{v} &= 0 \tag*{\text{\small[COM]}}\\
\partial_j v_j &= 0 \nonumber
\end{align}

\noindent\makebox[\linewidth]{\rule{\textwidth}{1pt}}

\[ \rho_0\left[\frac{\partial \underline{v}}{\partial t} + \nabla\cdot(\underline{v} \,\underline{v})\right] = \rho_0\underline{F} + \nabla\cdot\underline{\underline{T}}
\eqno\text{\small[COLM v1]}\]
\[ \; \rho_0 \left[\partial_t v_i + \partial_j(v_j v_i)\right] = \rho_0 F_i + \partial_j T_{ji} \]

\textbf{\emph{equivalently}} (using COM to simplify)
\[ \rho_0\frac{D \underline{v}}{D t} = \rho_0\underline{F} + \nabla\cdot\underline{\underline{T}} \eqno\text{\small[COLM v2]}\]
\[\rho_0 \left(\partial_t v_i + v_j\partial_j v_i\right) = \rho_0 F_i + \partial_j T_{ji}\]

\noindent\makebox[\linewidth]{\rule{\textwidth}{1pt}}

Don't be deceived because the last eq.\ looks just like its left counterpart. Here $\rho_0=const.$ (we're not solving for it), on the left $\rho = \rho(x,y,z,t)$ (we're solving for it).

\end{minipage}


\clearpage

\thispagestyle{fancy}


\noindent
The stress tensor of a \textbf{\emph{Newtonian fluid}} ($\undertilde{\text{terms}}$ vanish for \emph{incompressible} flows) in Cartesian coords $(x_1,x_2,x_3)=(x,y,z)$ with velocity $\underline{v} = \big(v_1(x,y,z),v_2(x,y,z),v_3(x,y,z)\big)$:
\begin{align}
\hspace{-5mm}
\underline{\underline{T}}
&=
\begin{pmatrix}
-p & 0 & 0\\
0 & -p & 0\\
0 & 0 & -p
\end{pmatrix} +
\begin{pmatrix}
2\mu \frac{\partial v_1}{\partial x} - \frac{2}{3}\mu\undertilde{\nabla\cdot\underline{v}} & \mu\big(\frac{\partial v_2}{\partial x} + \frac{\partial v_1}{\partial y}\big) & \mu\big(\frac{\partial v_3}{\partial x} + \frac{\partial v_1}{\partial z}\big)\\[3mm]
\mu\big(\frac{\partial v_1}{\partial y} + \frac{\partial v_2}{\partial x}\big) & 2\mu \frac{\partial v_2}{\partial y} - \frac{2}{3}\mu\undertilde{\nabla\cdot\underline{v}} & \mu\big(\frac{\partial v_3}{\partial y} + \frac{\partial v_2}{\partial z}\big)\\[3mm]
\mu\big(\frac{\partial v_1}{\partial z} + \frac{\partial v_3}{\partial x}\big) & \mu\big(\frac{\partial v_2}{\partial z} + \frac{\partial v_3}{\partial y}\big)  & 2\mu \frac{\partial v_3}{\partial z} - \frac{2}{3}\mu\undertilde{\nabla\cdot\underline{v}}
\end{pmatrix}
\nonumber\\[3mm]
T_{ij} &= -p\delta_{ij} - \tfrac{2}{3}\mu \partial_k v_k \delta_{ij} + 2\mu\partial_{(i}v_{j)}, \tag*{\text{\small[\underline{total} (Cauchy) stress tensor]}}
\end{align}
The Newtonian \emph{viscous} stress tensor is $\tau_{ij} = - \tfrac{2}{3}\mu \partial_k v_k \delta_{ij} + 2\mu\partial_{(i}v_{j)}$.
%\begin{align}
%\underline{\underline{\tau}}
%&=
%\begin{pmatrix}
%2\mu \frac{\partial v_1}{\partial x} - \frac{2}{3}\undertilde{\mu{\nabla\cdot\underline{v}}} & \mu\big(\frac{\partial v_1}{\partial y} + \frac{\partial v_2}{\partial x}\big) & \mu\big(\frac{\partial v_1}{\partial z} + \frac{\partial v_3}{\partial x}\big)\\[3mm]
%\mu\big(\frac{\partial v_2}{\partial x} + \frac{\partial v_1}{\partial y}\big) & 2\mu \frac{\partial v_2}{\partial y} - \frac{2}{3}\undertilde{\mu{\nabla\cdot\underline{v}}} & \mu\big(\frac{\partial v_2}{\partial z} + \frac{\partial v_3}{\partial y}\big)\\[3mm]
%\mu\big(\frac{\partial v_3}{\partial x} + \frac{\partial v_1}{\partial z}\big) & \mu\big(\frac{\partial v_3}{\partial y} + \frac{\partial v_2}{\partial z}\big)  & 2\mu \frac{\partial v_3}{\partial z} - \frac{2}{3}\mu\undertilde{\nabla\cdot\underline{v}}
%\end{pmatrix}
%\nonumber\\[3mm]
%\tau_{ij} &= - \tfrac{2}{3}\mu \partial_k v_k \delta_{ij} + 2\mu\partial_{(i}v_{j)} \tag*{\text{\small[viscous (Newtonian) stress tensor]}}
%\end{align}
\\
\underline{Note:} $\underline{\underline{\tau}}$ has \emph{nonzero} diagonal components, however, because of our \emph{definition} $p = -\frac{1}{3}\tr \underline{\underline{T}} = -\frac{1}{3}T_{ii}$ we must have $\tr\underline{\underline{\tau}} = \tau_{ii} = \tau_{11}+\tau_{22}+\tau_{33} = 0$
%Indeed,
%\[%\hspace{-1cm}
%\tau_{ii} %= 2\mu \frac{\partial u}{\partial x} - \frac{2}{3}\mu\nabla\cdot\underline{v} + 2\mu \frac{\partial v}{\partial y} - \frac{2}{3}\mu\nabla\cdot\underline{v} + 2\mu \frac{\partial w}{\partial z} - \frac{2}{3}\mu\nabla\cdot\underline{v}
%= 2\mu\left(\frac{\partial v_1}{\partial x}+\frac{\partial v_2}{\partial y}+\frac{\partial v_3}{\partial z} %\right) - 2\mu {\nabla\cdot\underline{v}} = 0 \eqno\text{\small[trace of viscous stress tensor]}
%\]
for \emph{both compressible and incompressible} flows, requiring the use of Stokes' hypothesis (that is, $\lambda = -2\mu/3$).
\\
\underline{Note:} $\rho$ and $\mu$ do \emph{not} have to be constant here, and this fluid could be compressible, in which case an \emph{equation of state} is needed to relate $p$ and $\rho$ (assuming isothermal conditions).


\vfill

For constant properties, $\rho=\rho_0$ and $\mu=\mu_0$, using the above $\underline{\underline{T}}$ in COLM v2, we obtain the \textbf{\emph{Navier--Stokes equations for an {incompressible flow of a Newtonian} fluid}}:
\begin{empheqboxed}
\begin{align}
&\nabla\cdot\underline{v} = 0,\qquad \rho_0 \frac{D \underline{v}}{D t} = \rho_0 \underline{F} -\nabla p + \mu_0 \nabla^2\underline{v}\qquad\qquad [\rho_0,\mu_0=const.] \tag*{\text{\small[iNS]}}\\[3mm]
&\partial_j v_j = 0, \hspace{38pt} \rho_0 \left(\partial_t v_i + v_j  \partial_j v_i \right) = \rho_0 F_i - \partial_i p + \mu_0 \partial_j \partial_j v_i \nonumber
\end{align}
\end{empheqboxed}

\end{document}
